% -----------------------------------------------------------------------------------------------------
\section{程函恒等式与 Wilson 线的 Ward 恒等式}\label{sec:appB-eikonal}

本节给出第~3~章\cref{lem:prelim-eikonal-identity}的证明。

\begin{lemma}[程函恒等式（洗牌积）]\label{lem:appB-shuffle}
(i) 设复数 $y_1,\dots,y_m$ 的任意非空子集之和都不为零，则
\begin{equation}\label{eq:appB-shuffle}
\sum_{\sigma\in S_m}\prod_{j=1}^m\frac{1}{y_{\sigma(1)}+\dots+y_{\sigma(j)}}=\prod_{j=1}^m\frac1{y_j}.
% VF-WL-4
\end{equation}
更一般地，对两个保持内部次序的序列 $(q_1,\dots,q_p)$ 与 $(r_1,\dots,r_q)$，记 $E(x_1,\dots,x_l)=\prod_{j=1}^l(x_1+\dots+x_j)^{-1}$，则对全部 $\binom{p+q}{p}$ 个洗牌 $\sigma$ 求和有
$\sum_\sigma E(\sigma)=E(q_1,\dots,q_p)\,E(r_1,\dots,r_q)$。
(ii) 取 $y_j=n\cdot k_j+\ii\epsilon$（$\epsilon>0$），则左边第 $j$ 个分母为 $n\cdot(k_{\sigma(1)}+\dots+k_{\sigma(j)})+\ii j\epsilon$，正是绝热正规子 (C3) 的累积形式：
程函恒等式 $\sum_\sigma\prod_j\frac{\ii}{n\cdot(k_{\sigma(1)}+\dots+k_{\sigma(j)})+\ii j\epsilon}=\prod_j\frac{\ii}{n\cdot k_j+\ii\epsilon}$ 在\emph{有限} $\epsilon$ 下严格成立（指向未来的线同理）。
(iii) 若每个程函分母都加\emph{同一个} $+\ii\epsilon$，恒等式只在 $\epsilon\to0$ 时按分布意义成立：$m=2$ 时（$a_j=n\cdot k_j$）
\[
\prod_j\frac{1}{a_j+\ii\epsilon}-\sum_\sigma\prod_j\frac{1}{a_{\sigma(1)}+\dots+a_{\sigma(j)}+\ii\epsilon}=\frac{-\ii\epsilon}{(a_1+\ii\epsilon)(a_2+\ii\epsilon)(a_1+a_2+\ii\epsilon)},
\]
它与 $\frac1{(1+a_1^2)(1+a_2^2)}$ 的配对为 $\frac{\pi^2\epsilon}{(1+\epsilon)^2(2+\epsilon)}$，只在 $\epsilon\to0$ 时为零。
% uses: none
\end{lemma}

\begin{proof}
(i) 对 $m$ 归纳。按最后一个元素 $\sigma(m)=l$ 分组：末因子 $1/(y_1+\dots+y_m)$ 公共，其余由归纳假设等于 $\prod_{j\ne l}1/y_j$，故左边
$=\frac{1}{\sum_jy_j}\sum_l\prod_{j\ne l}\frac1{y_j}=\prod_j\frac1{y_j}$。洗牌积同样对 $p+q$ 归纳（最后一个元素是 $q_p$ 或 $r_q$）。等价地，$\prod_j\frac{\ii}{y_j}=\prod_j\int_{-\infty}^0\dd\lambda_j\ee^{-\ii\lambda_jy_j}$
（$\Imag y_j>0$），而 $(-\infty,0]^m$ 是 $m!$ 个单形的不交并，每个单形给出左边一项：阿贝尔 Wilson 线是普通指数，$W^{(m)}=(W^{(1)})^m/m!$。
(ii) 由 (i)：$y_{\sigma(1)}+\dots+y_{\sigma(j)}$ 的虚部为 $j\epsilon$。(iii) 通分即得；配对用留数：对 $a_1$ 闭合于上半平面只有 $a_1=\ii$，给出 $\frac{-\ii\pi}{(1+\epsilon)(a_2+\ii(1+\epsilon))}$，
再对 $a_2$ 同样得 $\frac{-\pi}{(2+\epsilon)(1+\epsilon)}$，乘以 $-\ii\epsilon$ 即得（亦经数值核对）。% check: writing/work/W05/s4_wilson.py (e)
\end{proof}

即：只要每个累积分母的正规子等于已吸收胶子的正规子之和，恒等式在有限正规子下就严格成立——场的绝热阻尼正给出这种累积形式；朴素的统一 $+\ii\epsilon$（$\delta$ 正规化）不满足这一条件，恒等式随之失效（(iii)；$m=2$ 时这一条件也是必要的）
（参见\cite[式~(5)--(6)]{Echevarria:2015byo}；阿贝尔恒等式本身是经典的）。

\begin{remark}[半无限 Wilson 线的 Ward 恒等式]\label{rem:appB-ward}
在 $W[0,-\infty n]$ 的 $m$ 胶子项中（色矩阵 $T_1,\dots,T_m$ 已含顶点因子，$T_1$ 最靠近 $-\infty$；记该项为 $Y$），把一个动量为 $q$、极化为 $q^\mu$、色为 $T_q$ 的标量极化胶子插入全部 $m+1$ 个位置并求和，得
\begin{equation}\label{eq:appB-ward}
\sum_{p=0}^{m}(\text{插在第 }p\text{ 个胶子之后})=-g\,T_q\,Y+g\sum_{p=0}^{m-1}Z_p\;T_m\cdots T_{p+2}\,[T_q,T_{p+1}]\,T_p\cdots T_1 ,
% VF-WL-13
\end{equation}
其中 $Z_p$ 是第 $p+1$ 个胶子携带 $k_{p+1}+q$ 时的程函传播子之积。证明：内部插入用程函 Ward 恒等式
$\frac{\ii}{n\cdot(\ell+q)}(\ii g\,n\cdot q)\frac{\ii}{n\cdot\ell}=g[\frac{\ii}{n\cdot(\ell+q)}-\frac{\ii}{n\cdot\ell}]$ 拆开，最靠近 $-\infty$ 的插入给出 $-g$，相邻位置的项逐对相消，
只剩相邻色序之差（对易子项，与经三胶子顶点的图相消，附录~F）与端点项；阿贝尔情形和为 $-gY$。若标量极化也带同样的绝热移动，式~\eqref{eq:appB-ward} 在有限 $\epsilon$ 下严格成立（以 SU(3) 生成元数值核对至 $m=4$）。% check: s4_wilson.py (f)
\uses{prop:appB-rules}
\end{remark}

\begin{proposition}[非阿贝尔指数化 \tagA]\label{prop:appB-webs}
在一点相遇的 Wilson 线的真空关联函数（及其在背景场中的期望值）可以写成 web 之和的指数：$\langle W\rangle=\exp\big[\sum_w\tilde C(w)\,\mathcal F(w)\big]$，其中 $\mathcal F(w)$ 是 web 的运动学积分，
$\tilde C(w)$ 是修正的色因子；对两条线，web 是交换胶子在两条线上的连接不能被分成两个独立子集的图，其修正色因子是“最大非阿贝尔”部分\cite{Gatheral1983,FrenkelTaylor1984}。
\uses{prop:appB-rules,lem:appB-shuffle}
\end{proposition}

\begin{proofsketch}
阿贝尔情形：由\cref{prop:appB-rules}与洗牌积（\cref{lem:appB-shuffle}），独立的连接组之间因子化，$\langle W\rangle$ 是连通光子子图之和的指数。非阿贝尔情形：把每个图的色因子按交换次序写成对称部分与对易子之和，
对称部分用洗牌积重组进指数，对易子部分留下修正色因子 $\tilde C(w)$（两条线时为“最大非阿贝尔”部分）；$O(g^2)$ 时即\cref{prop:sudakov-webs}(b) 的恒等式（指针），两圈交叉梯形图的修正色因子为 $-\frac12C_FC_A$。
完整论证见\cite{Gatheral1983,FrenkelTaylor1984}。
\end{proofsketch}

\paragraph{地位与缺口}全阶的组合学论证（[A]），对两条线在文献中给出\cite{Gatheral1983,FrenkelTaylor1984}；在 $O(g^2)$ 是严格的代数恒等式（\cref{prop:sudakov-webs}(b)，[R, $R_{\mathrm{def}}$]）。
多条线的 web 及其重整化本书不需要。所需的有限正规子必须是绝热型的（\cref{lem:appB-shuffle}(iii)）。
